\documentclass[11pt,a4paper]{article}

% -------- packages --------
\usepackage[utf8]{inputenc}
\usepackage[T1]{fontenc}
\usepackage{lmodern}
\usepackage[margin=2.5cm]{geometry}
\usepackage{amsmath,amssymb,amsthm}
\usepackage{bm}
\usepackage{mathtools}
\usepackage{physics}
\usepackage{siunitx}
\usepackage{xcolor}
\usepackage{hyperref}
\usepackage{cleveref}
\usepackage{enumitem}
\usepackage{listings}
\usepackage{booktabs}
\usepackage{microtype}

\hypersetup{
    colorlinks=true,
    linkcolor=blue!60!black,
    citecolor=blue!60!black,
    urlcolor=blue!60!black
}

% -------- code listing style --------
\definecolor{codegray}{rgb}{0.5,0.5,0.5}
\definecolor{codebg}{rgb}{0.97,0.97,0.97}
\definecolor{codekw}{rgb}{0.13,0.29,0.53}
\definecolor{codestr}{rgb}{0.58,0.0,0.0}
\definecolor{codecom}{rgb}{0.25,0.50,0.35}

\lstdefinestyle{pythonstyle}{
    backgroundcolor=\color{codebg},
    basicstyle=\ttfamily\footnotesize,
    keywordstyle=\color{codekw}\bfseries,
    commentstyle=\color{codecom}\itshape,
    stringstyle=\color{codestr},
    numberstyle=\tiny\color{codegray},
    numbers=left,
    numbersep=8pt,
    frame=single,
    framesep=5pt,
    rulecolor=\color{codegray!40},
    breaklines=true,
    showstringspaces=false,
    language=Python,
    tabsize=4,
    captionpos=b
}
\lstset{style=pythonstyle}

% -------- macros --------
\newcommand{\G}{\bm{G}}
\newcommand{\T}{\bm{T}}
\newcommand{\Gam}{\bm{\Gamma}}
\newcommand{\dL}{\delta\bm{L}}
\newcommand{\uinc}{\bm{u}^{\text{inc}}}
\newcommand{\avg}[1]{\left\langle #1 \right\rangle}
\newcommand{\Oh}{O_h}
\DeclareMathOperator{\Im}{Im}
\DeclareMathOperator{\Re}{Re}

% -------- title --------
\title{Point-Scattering vs.\ Volume-Averaged Self-Consistent T-Matrix \\
       for Cubic-Voxel Elastodynamic Scattering \\
       \large Research notes and implementation roadmap}
\author{Tod Nestor}
\date{\today}

\begin{document}
\maketitle

\begin{abstract}
\noindent
Notes on the comparison between point-scatterer and volume-averaged
self-consistent T-matrix formulations for cubic voxels superimposed
on a homogeneous elastic medium, and a concrete roadmap for producing
the comparison within the existing \texttt{Thesis\_Version2} codebase.
The volumetric (mean-field) form is already implemented in
\texttt{effective\_contrasts.py}; the point-scatterer limit requires
only a minor modification that reuses the existing Gauss--Legendre
self-Green's tensor machinery. The definitive published reference is
Leroy, Chastanet \& Derode~\cite{LeroyChastanetDerode2010}; the
closest convergence-analysis precedent is the EM discrete-dipole
approximation literature (Yurkin \& Hoekstra~\cite{YurkinHoekstra2006a,
YurkinHoekstra2006b}), which has no published elastodynamic
counterpart.
\end{abstract}

\tableofcontents
\vspace{1em}
\hrule
\vspace{1em}

% =====================================================
\section{The direct elastodynamic match}
% =====================================================

The closest published comparison is Leroy, Chastanet \&
Derode~\cite{LeroyChastanetDerode2010}, \emph{``Mean-field T-matrix
approach to elastic wave scattering by small and point-like objects''}
(\emph{Waves in Random and Complex Media} \textbf{20}(4), 591--613,
2010). They derive the vector-elastic T-matrix under a volumetric
average of stress and momentum inside the inclusion (precisely the
mean-field / Eshelby-style dressing used in \texttt{effective\_contrasts.py}),
explicitly work out the point-scatterer limit of that same T-matrix,
and benchmark both against the exact Mie solution for elastic spheres.
Both forms satisfy the elastodynamic optical theorem, and the paper
comments specifically on low-frequency resonances being captured only
by the volumetric version for high-contrast inclusions --- the point
limit misses them. The EM precursor is
Pellegrini, Stout \& Thibaudeau~\cite{PellegriniStoutThibaudeau1997},
\emph{``Off-shell mean-field electromagnetic T-matrix of finite-size
spheres and fuzzy scatterers''}, on which Leroy et al.\ build the
off-shell formalism.

The canonical reference on the pathologies of the point-scatterer side
(radiative correction, renormalization, divergences) is
\emph{``Point scatterers for classical waves''}
by de Vries, van Coevorden \& Lagendijk~\cite{deVries1998}. The paper
is scalar/EM rather than elastic, but the regularization machinery
maps over directly and it is the standard reference for why naive
point scatterers violate the optical theorem without the imaginary
self-term.

% =====================================================
\section{The convergence story lives in the EM-DDA literature}
% =====================================================

No elastodynamic paper appears to exist that performs the full
Yurkin-style convergence study on cubic voxels. The EM-DDA work maps
onto the problem almost term-for-term, because the question
``how does the intervoxel propagator converge as $d \to 0$ under
different single-site dressings?'' is exactly the DDA polarizability
question.

\begin{itemize}[leftmargin=1.5em]
  \item Yurkin \& Hoekstra (2006)~\cite{YurkinHoekstra2006a,YurkinHoekstra2006b}
        --- rigorous convergence analysis (theory + extrapolation) in
        two parts. The key practical result: for cubically shaped
        scatterers (i.e.\ targets that exactly tile with cubic
        subvolumes, which is precisely the lattice-of-cubes situation),
        the linear-in-$d$ error term is strongly suppressed, and
        extrapolation gives $\sim\!10^2\times$ error reduction.
  \item Collinge \& Draine (2004)~\cite{CollingeDraine2004} ---
        \emph{``DDA with polarizabilities that account for both finite
        wavelength and target geometry''}. The FCD (filtered coupled
        dipole) prescription: the single-site dressing includes a
        finite-volume correction, and convergence is dominated by
        that choice, not the lattice.
  \item Piller (1998)~\cite{Piller1998} --- coupled-dipole for
        high-permittivity (equivalent to the high-contrast regime);
        a filtered Green's tensor regularizes the intervoxel coupling
        specifically to handle slow convergence for large contrast.
        This is the paper most relevant to the ``poor convergence if
        calculated the wrong way'' observation.
  \item Yurkin, Maltsev \& Hoekstra~\cite{YurkinReview2007} ---
        \emph{``The discrete dipole approximation: an overview and
        recent developments''}. Reviews LDR / CM / FCD / IGT
        polarizability options side-by-side with convergence curves;
        the single best starting point for seeing the trade-offs in
        one place.
\end{itemize}

\noindent
The IGT (Integrated Green Tensor) option is the EM analogue of the
Gauss--Legendre regularization currently in use --- so the Yurkin
papers directly assess the chosen method against alternatives.

% =====================================================
\section{Elastodynamic volume-integral papers worth citing}
% =====================================================

\begin{itemize}[leftmargin=1.5em]
  \item Gubernatis, Domany, Krumhansl \& Huberman
        (1977)~\cite{Gubernatis1977} --- Born approximation for
        elastic flaws, benchmarked against Mie for spheres. The
        companion paper~\cite{Gubernatis1977a} is the integral-equation
        foundation.
  \item Yvonnet (2011)~\cite{Yvonnet2011} --- ``A fast numerical
        strategy for computing the elastic properties of heterogeneous
        media based on voxel-based digital images: a
        space-Lippmann--Schwinger scheme''. Elastostatic voxel LS
        with per-voxel eigenstrain integration of the Green's tensor
        (not FFT). Directly relevant methodology for the self-term,
        awaiting a frequency-domain extension.
  \item Bonnet (2016)~\cite{Bonnet2016} --- ``Solvability of a volume
        integral equation formulation for anisotropic elastodynamic
        scattering''. Rigorous functional-analytic foundation.
  \item Korneev \& Johnson (1993, 1996)~\cite{KorneevJohnson1993,
        KorneevJohnson1996} --- exact P/S scattering from an elastic
        sphere. Standard validation target for any self-consistent
        T-matrix code. Useful as a regression test for cubic results
        in the sphere-like limit ($n_{\text{subcell}} \to \infty$
        on a spherical mask).
\end{itemize}

% =====================================================
\section{What does not appear to exist}
% =====================================================

A direct elastodynamic analogue of Yurkin--Hoekstra --- i.e.\ a
published convergence study for a cubic-voxel Lippmann--Schwinger /
Foldy--Lax elastic solver that sweeps point-scatterer vs.\
volume-averaged self-term across frequency and contrast, with error
bounds against an exact reference (sphere Mie via Korneev--Johnson).
The \texttt{Thesis\_Version2} infrastructure is arguably closer to
producing this result than anything currently in the literature. A
clean figure of ``relative error vs.\ $ka$ for (point / volumetric /
IGT-equivalent) self-terms, benchmarked against Korneev--Johnson for
a spherical inclusion tiled by cubes'' would be a genuine contribution
to the seismic-scattering literature. Leroy et al.\ did it only for
spheres, not for cubic voxels.

% =====================================================
\section{The mean-field T-matrix in code-matching form}
% =====================================================

Start from the elastodynamic Lippmann--Schwinger equation for a
single inclusion of volume $V$ and constant contrast $\dL$ (encoding
$\delta\rho$, $\delta\lambda$, $\delta\mu$), embedded in a homogeneous
background with Green's tensor $\G_0$:
\begin{equation}
  \bm{u}(\bm{x}) = \uinc(\bm{x})
    + \int_V \G_0(\bm{x} - \bm{x}') \, \dL \, \bm{u}(\bm{x}') \, dV'.
  \label{eq:LS}
\end{equation}

The mean-field approximation replaces the internal field
$\bm{u}(\bm{x}')$ inside $V$ by its volume average $\avg{\bm{u}}_V$
(plus, in the full vector version, the volume averages of strain and
rotation separately --- this is exactly the split that produces the
four amplification factors in the current implementation). Applying
$(1/V)\int_V \,dV$ to both sides of \cref{eq:LS} yields the closed
algebraic relation
\begin{equation}
  \avg{\bm{u}}_V = \avg{\uinc}_V + \Gam_0 \, \dL \, \avg{\bm{u}}_V,
\end{equation}
where the \emph{self-tensor} is the Green's tensor integrated twice
over the inclusion:
\begin{equation}
  \Gam_0 = \frac{1}{V} \int_V \int_V \G_0(\bm{x} - \bm{x}') \, dV' \, dV.
  \label{eq:Gamma0}
\end{equation}
Solving for $\avg{\bm{u}}_V$ and pushing back through the radiated
field gives the mean-field T-matrix:
\begin{equation}
  \boxed{
    \T_{\text{vol}}
    = V \, \dL \, \big[\, \bm{I} - \Gam_0 \, \dL \,\big]^{-1}.
  }
  \label{eq:Tvol}
\end{equation}

This is precisely the structure implemented in
\texttt{effective\_contrasts.py}. The four distinct amplification
factors (displacement, rotation, off-diagonal strain, diagonal strain)
come from decomposing $\avg{\bm{u}}_V$ and $\avg{\bm{\varepsilon}}_V$
into irreducible representations of the cube's $\Oh$ symmetry group.
A sphere supports two amplification factors
($\text{scalar} + \text{traceless-symmetric}$); on a cube the
traceless-symmetric representation further splits into diagonal
($E_g$) and off-diagonal ($T_{2g}$) pieces, giving
$T_1^c,\, T_2^c,\, T_3^c$. The implementation therefore has ``more
structure'' than the sphere version of Leroy--Derode --- it is the
correct generalization for $\Oh$ rather than $SO(3)$.

% =====================================================
\section{The point limit}
% =====================================================

The point-scatterer T-matrix is obtained by taking $V \to 0$ while
keeping a \emph{renormalized} strength $\bm{\tau} = V\,\dL$ finite. In
this limit the real part of $\Gam_0$ diverges --- the Green's tensor
is $O(1/r)$ in 3D, so $\int\!\!\int \G_0 \sim V^{2/3}$ times a
logarithmic-like correction that requires regularization --- and the
paper performs the renormalization by keeping only the imaginary part
of $\Gam_0$, which is finite and equals the radiative reaction
required by the optical theorem:
\begin{equation}
  \T_{\text{point}} = \bm{\tau}
    \, \big[\, \bm{I} - i\,\Im\{\Gam_0^{\text{point}}\} \, \bm{\tau} \,\big]^{-1}.
  \label{eq:Tpoint}
\end{equation}
Here $\Im\{\Gam_0^{\text{point}}\}$ is the combined P-wave and S-wave
radiation-damping tensor, schematically
\begin{equation}
  \Im\!\left\{\Gamma_{0,ij}^{\text{point}}\right\}
    = \frac{1}{6\pi \rho_0 \omega^2} \left(
        \frac{k_P^3}{3}\,\delta_{ij}
      + \frac{2 k_S^3}{3}\,\delta_{ij}^{(\text{pol-weighted})}
    \right),
\end{equation}
with the full tensor decomposition in eqs.\ (12)--(14) of
\cite{LeroyChastanetDerode2010}; the de~Vries--van
Coevorden--Lagendijk derivation~\cite{deVries1998}
(eqs.\ 3.5--3.11) makes clear why this structure is forced by the
optical theorem.

% =====================================================
\section{A minimal modification to \texttt{effective\_contrasts.py}}
% =====================================================

The existing $\Gam_0$ routine (Gauss--Legendre over the cube) returns
a complex tensor whose real part comes from the principal-value
integral of the singular piece of $\G_0$, and whose imaginary part
comes from the radiating piece $\sim e^{ikr}/r$. The point-limit
version zeros the real part and replaces $V\dL$ with the renormalized
$\bm{\tau}$:

\begin{lstlisting}[caption={Point-scatterer limit: drop-in companion to \texttt{compute\_cube\_tmatrix}.}]
def compute_cube_tmatrix_point_limit(omega, rho0, lam0, mu0,
                                      dlam, dmu, drho, a,
                                      renormalize="leroy"):
    """Point-scatterer limit of the mean-field cube T-matrix.

    Matches compute_cube_tmatrix signature; returns the same
    dataclass, with amplification factors computed in the
    V -> 0, tau = V * dL finite limit.
    """
    V  = (2.0 * a) ** 3
    dL = build_contrast_operator(dlam, dmu, drho, omega)

    # Full (volumetric) version would use Gamma0 = Re + i*Im.
    # Point limit drops the Re part (principal-value),
    # keeping only the radiation term.
    Gamma0_full  = gauss_legendre_self_greens(omega, rho0, lam0, mu0, a)
    Gamma0_point = 1j * Gamma0_full.imag                 # radiation only

    tau = V * dL                                         # renormalized strength
    I   = np.eye(tau.shape[0])
    T_point = tau @ np.linalg.inv(I - Gamma0_point @ tau)

    return package_as_cube_tmatrix_result(T_point, kind="point_limit")
\end{lstlisting}

\subsection{Two subtleties worth flagging}

\paragraph{Which \texorpdfstring{$\Im$}{Im} to keep.}
In Leroy--Derode, the radiation reaction in $\Gam_0$ has separate P-
and S-wave contributions with different $k^3$ scalings. The existing
Gauss--Legendre integrator computes these automatically as the
imaginary part of the integrand, so the one-liner
\texttt{1j * Gamma0.imag} is correct --- but only if
\texttt{gauss\_legendre\_self\_greens} integrates the \emph{radiating}
Green's tensor $e^{ikr}/(4\pi r)\cdot(\ldots)$ and not just the static
$1/r$ piece. Worth a unit test: in the limit $\omega \to 0$,
\texttt{Gamma0\_point} should vanish, while \texttt{Gamma0\_full}
should approach the Eshelby tensor for a cube (times $1/\mu_0$ with
the standard combinations).

\paragraph{The renormalization choice.}
``Point limit'' is not unique. Leroy--Derode regularize by retaining
the cube's volume in $\bm{\tau} = V\dL$, but other papers
renormalize through a scattering length $f_s$. Matching the Foldy
point-scatterer convention exactly requires $\bm{\tau}$ to be
replaced by $-4\pi f_s / k^2$ with $f_s$ extracted from the Rayleigh
limit of the full cube T-matrix. The two conventions give identical
far-field scattering but different interpretations of the
``strength'' parameter --- which matters when comparing resonance
frequencies.

% =====================================================
\section{The comparison plot}
% =====================================================

With both variants exposed, the definitive figure is:
\begin{quote}
\textbf{Relative error of the coherent P-wave far-field amplitude, as
a function of $k_P a$, for:}
\begin{itemize}[leftmargin=1.5em,nosep]
  \item Volumetric T-matrix (existing \texttt{compute\_cube\_tmatrix})
  \item Point-limit T-matrix (the version above)
  \item Born approximation (trivial limit, $\T = V\dL$)
        --- as a sanity floor
  \item \textbf{Reference:} Korneev--Johnson Mie for a sphere of
        equivalent volume $V_{\text{sphere}} = V_{\text{cube}}$
\end{itemize}
\end{quote}

Three sweeps across contrast regimes:
\begin{center}
\begin{tabular}{lc}
  \toprule
  Regime    & $|\delta\mu/\mu_0|$ \\
  \midrule
  Weak      & $\sim 0.1$ \\
  Moderate  & $\sim 1$   \\
  Strong    & $\sim 10$ (resonances appear) \\
  \bottomrule
\end{tabular}
\end{center}

\subsection{Expected qualitative result}

From the Leroy--Derode paper, on a sphere benchmark: the point limit
is accurate for weak contrast at all $ka < 0.3$, but fails
catastrophically at strong contrast where the volumetric version
captures the monopole resonance correctly. The cube case should
show the same split --- with the additional diagnostic that the
cube's $T_1^c / T_2^c / T_3^c$ decomposition reveals the
$\Oh$-point-group-induced resonance splitting that a sphere cannot
exhibit. Once this figure is produced, the thesis has both
\begin{enumerate}[nosep]
  \item validation of the volumetric approach on a known reference,
  \item a novel result --- the cube-vs-sphere resonance structure ---
        that is not in the literature.
\end{enumerate}

% =====================================================
\section{Implementation checklist}
% =====================================================

\begin{enumerate}[leftmargin=1.8em]
  \item Add \texttt{compute\_cube\_tmatrix\_point\_limit} as a
        drop-in companion in \texttt{effective\_contrasts.py}.
  \item Unit test: $\omega \to 0$ limit of the point version
        $\to 0$; volumetric version $\to$ cube Eshelby tensor
        (compare with known closed-form Eshelby for a cube if using
        one; otherwise converge against a fine FEM static reference).
  \item Unit test: optical theorem for both variants at a sweep of
        $ka$ values. Both should satisfy it to quadrature precision.
  \item Benchmark: spherical-mask cube lattice of volume
        $V_{\text{sphere}}$, compare against
        Korneev--Johnson~\cite{KorneevJohnson1993}.
  \item Produce error-vs-$ka$ plot at three contrast regimes.
  \item Reuse \texttt{scattering\_order\_decomposition} to generate
        Neumann-series convergence curves for both T-matrix variants
        --- these expose the difference in coupling strength between
        the two formulations.
\end{enumerate}

% =====================================================
\begin{thebibliography}{99}

\bibitem{LeroyChastanetDerode2010}
V.~Leroy, A.~Chastanet, and A.~Derode,
\emph{Mean-field T-matrix approach to elastic wave scattering by
small and point-like objects},
Waves in Random and Complex Media \textbf{20}(4), 591--613 (2010).

\bibitem{PellegriniStoutThibaudeau1997}
Y.~Pellegrini, D.~Stout, and P.~Thibaudeau,
\emph{Off-shell mean-field electromagnetic T-matrix of finite-size
spheres and fuzzy scatterers},
J. Phys.: Condens. Matter \textbf{9}, 177--191 (1997).

\bibitem{deVries1998}
P.~de Vries, D.~V.~van Coevorden, and A.~Lagendijk,
\emph{Point scatterers for classical waves},
Rev. Mod. Phys. \textbf{70}, 447--466 (1998).

\bibitem{YurkinHoekstra2006a}
M.~A.~Yurkin and A.~G.~Hoekstra,
\emph{Convergence of the discrete dipole approximation. I.\
Theoretical analysis},
J. Opt. Soc. Am. A \textbf{23}(10), 2578--2591 (2006);
\href{https://arxiv.org/abs/0704.0033}{arXiv:0704.0033}.

\bibitem{YurkinHoekstra2006b}
M.~A.~Yurkin, V.~P.~Maltsev, and A.~G.~Hoekstra,
\emph{Convergence of the discrete dipole approximation. II.\
An extrapolation technique to increase the accuracy},
J. Opt. Soc. Am. A \textbf{23}(10), 2592--2601 (2006);
\href{https://arxiv.org/abs/0704.0035}{arXiv:0704.0035}.

\bibitem{YurkinReview2007}
M.~A.~Yurkin and A.~G.~Hoekstra,
\emph{The discrete dipole approximation: an overview and recent
developments},
J. Quant. Spectrosc. Radiat. Transf. \textbf{106}, 558--589 (2007);
\href{https://arxiv.org/abs/0704.0038}{arXiv:0704.0038}.

\bibitem{CollingeDraine2004}
M.~J.~Collinge and B.~T.~Draine,
\emph{Discrete-dipole approximation with polarizabilities that account
for both finite wavelength and target geometry},
J. Opt. Soc. Am. A \textbf{21}, 2023--2028 (2004).

\bibitem{Piller1998}
N.~B.~Piller,
\emph{Coupled-dipole approximation for high permittivity materials},
Opt. Commun. \textbf{160}, 10--14 (1998).

\bibitem{Gubernatis1977}
J.~E.~Gubernatis, E.~Domany, J.~A.~Krumhansl, and M.~Huberman,
\emph{The Born approximation in the theory of the scattering of
elastic waves by flaws},
J. Appl. Phys. \textbf{48}, 2812--2819 (1977).

\bibitem{Gubernatis1977a}
J.~E.~Gubernatis, E.~Domany, and J.~A.~Krumhansl,
\emph{Formal aspects of the theory of the scattering of ultrasound
by flaws in elastic materials},
J. Appl. Phys. \textbf{48}, 2804--2811 (1977).

\bibitem{Yvonnet2011}
J.~Yvonnet,
\emph{A fast numerical strategy for computing the elastic properties
of heterogeneous media based on voxel-based digital images: a
space-Lippmann--Schwinger scheme},
Int. J. Numer. Methods Eng. \textbf{86}, 1--12 (2011).

\bibitem{Bonnet2016}
M.~Bonnet,
\emph{Solvability of a volume integral equation formulation for
anisotropic elastodynamic scattering},
J. Integral Equ. Appl. \textbf{28}(2), 169--203 (2016).

\bibitem{KorneevJohnson1993}
V.~A.~Korneev and L.~R.~Johnson,
\emph{Scattering of elastic waves by a spherical inclusion --- I.\
Theory and numerical results},
Geophys. J. Int. \textbf{115}, 230--250 (1993).

\bibitem{KorneevJohnson1996}
V.~A.~Korneev and L.~R.~Johnson,
\emph{Scattering of P and S waves by a spherically symmetric inclusion},
Pure Appl. Geophys. \textbf{147}, 675--718 (1996).

\end{thebibliography}

\end{document}
